% !Mode:: "TeX:UTF-8"
% 定理环境专项测试。42 变体只覆盖 theorem/definition/axiom/lemma/remark 五个，
% proof 一个都没有，而 proof 恰好是换 thmtools 时差别最大的地方（ntheorem 的
% thmmarks 与 amsthm 的 \qedhere 是两套机制）。换引擎前后拿这份文档逐字节比。
\documentclass[fontset=fandol,type=master,campus=harbin]{hithesisbook}
\hitsetup{info={ctitle={定理测试},cauthor={测试},csupervisor={测试},
  cxueke={工学},csubject={测试},caffil={测试},ckeywords={甲,乙},
  etitle={Theorem Test},eauthor={T},esupervisor={T},exueke={E},
  esubject={E},eaffil={E},ekeywords={a,b}}}
\newcommand\filler{这是用来把定理撑到跨页的填充文字。}
\begin{document}
\begin{cabstract}摘要。\end{cabstract}
\begin{eabstract}Abstract.\end{eabstract}
\makecover
\tableofcontents
\mainmatter
\chapter{全部环境}

\section{十三个编号环境，不带注记}
\begin{assumption}\label{thm:as}假设正文。\end{assumption}
\begin{definition}定义正文。\end{definition}
\begin{proposition}命题正文。\end{proposition}
\begin{lemma}引理正文。\end{lemma}
\begin{theorem}定理正文。\end{theorem}
\begin{axiom}公理正文。\end{axiom}
\begin{corollary}推论正文。\end{corollary}
\begin{exercise}练习正文。\end{exercise}
\begin{example}例正文。\end{example}
\begin{remark}注释正文。\end{remark}
\begin{problem}问题正文。\end{problem}
\begin{conjecture}猜想正文。\end{conjecture}
\begin{fact}事实正文。\end{fact}

\section{带注记}
\begin{theorem}[（霍金）]带中文注记。\end{theorem}
\begin{definition}[Hawking]带英文注记。\end{definition}
\begin{lemma}[$\alpha>0$]带数学注记。\end{lemma}

\section{交叉引用}
见假设~\ref{thm:as}、\autoref{thm:as}。

\section{proof 的四种收尾}
\begin{proof}以文字结尾。\end{proof}

\begin{proof}以行内公式结尾 $a+b=c$。\end{proof}

\begin{proof}
以显示公式结尾：
\begin{equation}
  E=mc^2
\end{equation}
\end{proof}

\begin{proof}
以列表结尾：
\begin{itemize}
  \item 甲
  \item 乙
\end{itemize}
\end{proof}

\section{嵌套与跨页}
\begin{theorem}
外层定理。
\begin{proof}内层证明。\end{proof}
\end{theorem}

\begin{theorem}
跨页定理：
\filler\filler\filler\filler\filler\filler\filler\filler\filler\filler
\filler\filler\filler\filler\filler\filler\filler\filler\filler\filler
\filler\filler\filler\filler\filler\filler\filler\filler\filler\filler
\filler\filler\filler\filler\filler\filler\filler\filler\filler\filler
\end{theorem}

\begin{proof}
跨页证明：\filler\filler\filler\filler\filler\filler\filler\filler
\filler\filler\filler\filler\filler\filler\filler\filler\filler\filler
\filler\filler\filler\filler\filler\filler\filler\filler\filler\filler
\filler\filler\filler\filler\filler\filler\filler\filler\filler\filler
\end{proof}

\section{自定义环境}
\newtheorem{nonsense}{胡说}[chapter]
\begin{nonsense}手册里教的写法。\end{nonsense}
\newtheorem{nonsensetwo}[nonsense]{胡说二}
\begin{nonsensetwo}共用计数器的写法。\end{nonsensetwo}
\newtheorem*{nonum}{无编号}
\begin{nonum}无编号环境。\end{nonum}
\end{document}
